% Бүлэг 3
\chapter{СИСТЕМИЙН ЗОХИОН БАЙГУУЛАЛТ БА ХЭРЭГЖҮҮЛЭЛТ}
\label{Chapter3}
\pagecolor{white}

%-------------------------------------------------------------------------------

\section{Системийн шаардлагын дүн шинжилгээ}

\subsection{Функциональ шаардлага}

Системийн функциональ шаардлагыг \textbf{Хэрэглэгч} (иргэн) болон \textbf{Админ} (заалны цаг олгогч хянагч) гэсэн хоёр үүрэгт тулгуурлан тодорхойлсон. Захиалгын удирдлага хэрэглэгчид болон автомат cron ажлаар явагддаг бөгөөд Админ нь \code{fixed\_bookings} цуглуулгаар байгууллагын тогтмол захиалгуудыг тохируулдаг.

\textbf{Хэрэглэгчийн функцүүд:}
\begin{enumerate}
    \item Email хаяг болон нууц үгээр бүртгүүлэх; бүртгэлийн үед тухайн email рүү 6 оронтой баталгаажуулах код илгээж, кодоо оруулсны дараа нэвтрэх боломж нээгдэх;
    \item Таван спорт заалны жагсаалтыг нүүр дэлгэцэд харах;
    \item Тухайн заалны боломжтой болон хаагдсан цагийн хуваарийг харах;
    \item Цаг захиалах (заал, огноо, цагийн үе сонгох), DBK-xxxxx форматын захиалгын код авах;
    \item Идэвхтэй захиалгаа цуцлах;
    \item Захиалгын түүхийг захиалгын дэлгэцэд харах;
    \item Захиалгаас 1 цагийн өмнө локал мэдэгдэл хүлээн авах;
    \item AI чатботтой спорт заалтай холбоотой асуулт асуух.
\end{enumerate}

\textbf{Админы функцүүд (заалны цаг олгогч хянагч):}
\begin{enumerate}
    \item Байгууллага, компаниудын долоо хоногийн тогтмол захиалгыг \code{fixed\_bookings} цуглуулгад нэмэх, устгах, өөрчлөх;
    \item Хаалттай цагийн үеийг тодорхойлж цагийн хуваарьт тусгах.
\end{enumerate}

\textbf{Системийн автомат функцүүд:}
\begin{enumerate}
    \item 15 минут тутамд cron ажил дуусгавар болсон захиалгыг \code{completed} төлөвт шилжүүлнэ;
    \item Байгууллагуудын тогтмол (давтагдах) захиалгыг \code{fixed\_bookings} цуглуулгаас уншиж долоо хоног бүр цагийн хуваарьт тусгана.
\end{enumerate}

\subsection{Функциональ бус шаардлага}

Sommerville (2019) \cite{software_engineering}-д тодорхойлсон аргачлалын дагуу функциональ бус шаардлагыг хэмжигдэхүйц хэлбэрээр тодорхойлсон:

\begin{table}[h]
\centering
\caption{Функциональ бус шаардлагуудын хүснэгт}
\label{tab:nonfunctional}
\begin{tabularx}{\textwidth}{|p{2.8cm}|X|p{3cm}|}
\hline
\textbf{Ангилал} & \textbf{Шаардлага} & \textbf{Хэмжих арга} \\
\hline
Гүйцэтгэл & Апп 3 секундын дотор ачаалах; API хариулт 2 секундын дотор & Chrome DevTools, хэрэглэгчийн цагийн хэмжилт \\
\hline
Найдвартай байдал & Firebase 99.9\%, AWS EC2 99.5\% uptime баталгаа & AWS CloudWatch \\
\hline
Аюулгүй байдал & Нэвтрэлт Firebase Auth-аар; API нууцлал GitHub Secrets + .env & OWASP Mobile Security \\
\hline
Хэрэглэхэд хялбар & Material Design 3 стандарт; 5/5-аас 4.0+ оноо & SUS (Usability Scale) \\
\hline
Дэмжих платформ & Android 8.0+ (API 26+) болон iOS 13.0+ & Бодит төхөөрөмж туршилт \\
\hline
Масштабчлал & AWS EC2 instance-г шаардлагаар томруулах боломж & AWS Auto Scaling \\
\hline
\end{tabularx}
\end{table}

%-------------------------------------------------------------------------------

\section{Системийн UML диаграммууд}

\subsection{Use Case Diagram (Хэрэглэгчийн тохиолдлын диаграмм)}

Use Case Diagram нь системийн гадаад оролцогч болон тэдгээрийн гүйцэтгэх үндсэн функцуудыг тодорхойлно. Энэхүү системд \textbf{Хэрэглэгч} (иргэн) болон \textbf{Админ} (заалны цаг олгогч хянагч) гэсэн хоёр үндсэн actor оролцоно.

\begin{figure}[h]
\centering
\begin{tikzpicture}[
  uc/.style={draw, ellipse, minimum width=4.4cm, minimum height=0.85cm,
             text centered, font=\small, fill=blue!6, draw=blue!40},
  ucadm/.style={draw, ellipse, minimum width=4.4cm, minimum height=0.85cm,
             text centered, font=\small, fill=orange!10, draw=orange!50},
  lbl/.style={font=\small\bfseries},
  inc/.style={dashed, -latex', gray, thin}
]

% System boundary
\draw[thick, rounded corners=10pt, fill=gray!3] (-0.3,-0.5) rectangle (10.5,12.2);
\node[lbl, font=\normalsize\bfseries] at (5.1,12.6) {ДЗ-Заал Цаг Захиалгын Систем};

% ======= Хэрэглэгч Actor =======
\draw (-2.2,9.5) circle (0.38cm);
\draw (-2.2,9.12) -- (-2.2,8.1);
\draw (-2.85,8.65) -- (-2.2,8.85) -- (-1.55,8.65);
\draw (-2.2,8.1) -- (-2.65,7.3);
\draw (-2.2,8.1) -- (-1.75,7.3);
\node[lbl, below, font=\small] at (-2.2,7.1) {Хэрэглэгч};

% ======= Админ Actor =======
\draw (12.5,9.5) circle (0.38cm);
\draw (12.5,9.12) -- (12.5,8.1);
\draw (11.85,8.65) -- (12.5,8.85) -- (13.15,8.65);
\draw (12.5,8.1) -- (12.05,7.3);
\draw (12.5,8.1) -- (12.95,7.3);
\node[lbl, below, font=\small, align=center] at (12.5,7.0) {Админ\\(Цаг хянагч)};

% ======= Хэрэглэгчийн use cases =======
\node[uc] (uc1) at (5,11.4) {Email бүртгэл / нэвтрэх};
\node[uc] (uc2) at (5,10.2) {Спорт заалнуудыг харах};
\node[uc] (uc3) at (5,9.0)  {Цагийн хуваарь харах};
\node[uc] (uc4) at (5,7.8)  {Цаг захиалах};
\node[uc] (uc5) at (5,6.6)  {Захиалга цуцлах};
\node[uc] (uc6) at (5,5.4)  {Захиалгын түүх харах};
\node[uc] (uc7) at (5,4.2)  {Локал мэдэгдэл авах};
\node[uc] (uc8) at (5,3.0)  {AI чатботтой харилцах};

% ======= Админы use cases =======
\node[ucadm] (adm1) at (5,1.8)  {Тогтмол захиалга удирдах};
\node[ucadm] (adm2) at (5,0.6)  {Хаалттай цаг тохируулах};

% ======= Хэрэглэгч холболт =======
\foreach \uc in {uc1,uc2,uc3,uc4,uc5,uc6,uc7,uc8}
  \draw (-1.82,8.8) -- (\uc.west);

% ======= Админ холболт =======
\foreach \adm in {adm1,adm2}
  \draw (12.12,8.8) -- (\adm.east);

% ======= include харилцаа =======
\draw[inc] (uc4.south) -- node[right,font=\scriptsize]{<<include>>} (uc3.south east);
\draw[inc] (uc1.east) -- ++(0.3,0) |- node[right,font=\scriptsize]{<<extend>>} (uc7.east);

\end{tikzpicture}
\caption{Use Case Diagram --- Хэрэглэгч болон Админы харилцаа}
\label{fig:usecase}
\end{figure}

%-------------------------------------------------------------------------------

\subsection{Цаг захиалах үйл ажиллагааны диаграмм (Activity Diagram)}

Цаг захиалах үйл явцын дарааллыг үйл ажиллагааны диаграммаар дүрсэлсэн. Хэрэглэгч нэвтэрсний дараа заал, огноо, цагийн үе сонгож захиалга үүсгэнэ. Сонгосон цаг аль хэдийн захиалагдсан бол систем алдаа мэдэгдэж өөр цаг сонгохыг хүснэ.

\begin{figure}[h]
\centering
\begin{tikzpicture}[
  proc/.style={draw, rounded corners=5pt, fill=yellow!10,
               minimum width=5cm, minimum height=0.8cm,
               text centered, font=\small},
  dec/.style={draw, diamond, fill=orange!15, aspect=2.5,
              text centered, font=\scriptsize, inner sep=1pt},
  term/.style={draw, circle, fill=black, minimum size=0.4cm},
  termend/.style={draw, circle, fill=black, minimum size=0.4cm,
                  double, double distance=2pt},
  arr/.style={-latex', thick}
]
\node[term]  (start)   at (0,12)   {};
\node[proc]  (login)   at (0,10.8) {Email/нууц үгээр нэвтрэх};
\node[proc]  (venue)   at (0,9.6)  {Спорт заал сонгох};
\node[proc]  (date)    at (0,8.4)  {Огноо сонгох};
\node[proc]  (slot)    at (0,7.2)  {Цагийн үе сонгох};
\node[dec]   (avail)   at (0,5.8)  {Цаг боломжтой?};
\node[proc]  (err)     at (4,5.8)  {Алдаа харуулах};
\node[proc]  (create)  at (0,4.4)  {Захиалга үүсгэх (Firestore)};
\node[proc]  (code)    at (0,3.2)  {DBK-xxxxx код авах};
\node[proc]  (notif)   at (0,2.0)  {1 цагийн өмнө мэдэгдэл товлох};
\node[proc]  (done)    at (0,0.8)  {Баталгаажуулалт харуулах};
\node[termend] (end)   at (0,-0.3) {};

\draw[arr] (start) -- (login);
\draw[arr] (login) -- (venue);
\draw[arr] (venue) -- (date);
\draw[arr] (date)  -- (slot);
\draw[arr] (slot)  -- (avail);
\draw[arr] (avail) -- node[right,font=\scriptsize]{Тийм} (create);
\draw[arr] (avail) -- node[above,font=\scriptsize]{Үгүй} (err);
\draw[arr] (err.north) |- (slot.east);
\draw[arr] (create) -- (code);
\draw[arr] (code)   -- (notif);
\draw[arr] (notif)  -- (done);
\draw[arr] (done)   -- (end);
\end{tikzpicture}
\caption{Цаг захиалах үйл ажиллагааны диаграмм}
\label{fig:activity}
\end{figure}

%-------------------------------------------------------------------------------

\subsection{Захиалга үүсгэх дарааллын диаграмм (Sequence Diagram)}

Захиалга үүсгэх үйл явцын оролцогч тал бүрийн харилцан үйлдлийн дарааллыг дарааллын диаграммаар илэрхийлсэн.

\begin{figure}[h]
\centering
\begin{tikzpicture}[
  lifeline/.style={draw, rectangle, minimum width=2.2cm, minimum height=0.7cm,
                   text centered, font=\small\bfseries, fill=blue!10},
  msg/.style={-latex', thick},
  ret/.style={-latex', dashed, thick},
  font=\small
]
% Lifeline headers
\node[lifeline] (U)  at (0,10)    {Хэрэглэгч};
\node[lifeline] (A)  at (3.5,10)  {Flutter App};
\node[lifeline] (F)  at (7,10)    {Firestore};
\node[lifeline] (B)  at (10.5,10) {Backend};

% Lifeline bodies (vertical dashed lines)
\draw[dashed] (0,9.65)   -- (0,-0.5);
\draw[dashed] (3.5,9.65) -- (3.5,-0.5);
\draw[dashed] (7,9.65)   -- (7,-0.5);
\draw[dashed] (10.5,9.65)-- (10.5,-0.5);

% Messages
\draw[msg] (0,9)    -- node[above,font=\scriptsize]{Захиалах дарах} (3.5,9);
\draw[msg] (3.5,8.2)-- node[above,font=\scriptsize]{Боломжтой цаг шалгах} (7,8.2);
\draw[ret] (7,7.5)  -- node[above,font=\scriptsize]{Хаагдсан цагийн жагсаалт} (3.5,7.5);
\draw[msg] (3.5,6.7)-- node[above,font=\scriptsize]{Боломжтой цаг харуулах} (0,6.7);
\draw[msg] (0,5.9)  -- node[above,font=\scriptsize]{Цаг батлах} (3.5,5.9);
\draw[msg] (3.5,5.1)-- node[above,font=\scriptsize]{bookings collection-д бичих} (7,5.1);
\draw[ret] (7,4.4)  -- node[above,font=\scriptsize]{DBK-xxxxx код} (3.5,4.4);
\draw[msg] (3.5,3.6)-- node[above,font=\scriptsize]{Мэдэгдэл товлох хүсэлт} (10.5,3.6);
\draw[ret] (10.5,2.9)--node[above,font=\scriptsize]{200 OK} (3.5,2.9);
\draw[ret] (3.5,2.1)-- node[above,font=\scriptsize]{Баталгаажуулалт + код харуулах} (0,2.1);
\draw[msg] (3.5,1.2)-- node[above,font=\scriptsize]{Local notification schedule} (3.5,0.4);
\node[font=\scriptsize,right] at (3.5,0.8) {(flutter\_local\_notifications)};
\end{tikzpicture}
\caption{Захиалга үүсгэх дарааллын диаграмм}
\label{fig:sequence}
\end{figure}

%-------------------------------------------------------------------------------

\subsection{Email баталгаажуулалтын дарааллын диаграмм}

Хэрэглэгч анх бүртгүүлэх үед тохиодог email баталгаажуулалтын дарааллыг дараах диаграммаар илэрхийлэв. Энэ үйл явц нь Firebase Auth, Node.js бэкэнд (Nodemailer) болон Firestore гурвыг хамардаг тул нарийвчилж харуулах шаардлагатай.

\begin{figure}[h]
\centering
\begin{tikzpicture}[
  lifeline/.style={draw, rectangle, minimum width=2.4cm, minimum height=0.7cm,
                   text centered, font=\small\bfseries, fill=orange!12},
  msg/.style={-latex', thick, blue!70},
  ret/.style={-latex', dashed, thick, green!50!black},
  err/.style={-latex', dashed, thick, red!60},
  font=\small
]
% Lifeline headers
\node[lifeline] (U)  at (0,11)    {Хэрэглэгч};
\node[lifeline] (A)  at (3,11)    {Flutter App};
\node[lifeline] (FB) at (6.5,11)  {Firebase Auth};
\node[lifeline] (B)  at (10,11)   {Node.js EC2};
\node[lifeline] (E)  at (13.5,11) {Gmail SMTP};

% Lifeline bodies
\foreach \x in {0, 3, 6.5, 10, 13.5}
  \draw[dashed, gray] (\x,10.65) -- (\x,-0.5);

% Step 1: Register
\draw[msg] (0,10.2)   -- node[above,font=\scriptsize]{Нэр, email, нууц үг} (3,10.2);
\draw[msg] (3,9.5)    -- node[above,font=\scriptsize]{createUserWithEmailAndPassword} (6.5,9.5);
\draw[ret] (6.5,8.8)  -- node[above,font=\scriptsize]{UID үүсгэгдсэн} (3,8.8);

% Step 2: Send code via backend
\draw[msg] (3,8.1)    -- node[above,font=\scriptsize]{POST /api/auth/send-code \{uid, email\}} (10,8.1);
\draw[msg] (10,7.3)   -- node[above,font=\scriptsize]{Nodemailer: 6 оронтой код} (13.5,7.3);
\draw[ret] (13.5,6.6) -- node[above,font=\scriptsize]{Илгээлт амжилттай} (10,6.6);

% Step 3: Code saved in Firestore
\draw[msg] (10,5.8)   -- node[above,font=\scriptsize]{email\_verifications-д код хадгалах} (6.5,5.8);
\draw[ret] (6.5,5.1)  -- node[above,font=\scriptsize]{200 OK} (3,5.1);
\draw[ret] (3,4.4)    -- node[above,font=\scriptsize]{Кодын оруулах дэлгэц харуулах} (0,4.4);

% Step 4: User enters code
\draw[msg] (0,3.7)    -- node[above,font=\scriptsize]{6 оронтой кодоо оруулах} (3,3.7);
\draw[msg] (3,2.9)    -- node[above,font=\scriptsize]{POST /api/auth/verify-code} (10,2.9);
\draw[ret] (10,2.2)   -- node[above,font=\scriptsize]{emailVerified:true $\rightarrow$ Firestore} (3,2.2);
\draw[ret] (3,1.4)    -- node[above,font=\scriptsize]{Нэвтрэх боломж нээгдсэн} (0,1.4);
\draw[msg] (10,0.6)   -- node[above,font=\scriptsize]{email\_verifications document устгах} (6.5,0.6);

\end{tikzpicture}
\caption{Email баталгаажуулалтын дарааллын диаграмм (Firebase Auth + Nodemailer)}
\label{fig:email_seq}
\end{figure}

%-------------------------------------------------------------------------------

\subsection{Firestore өгөгдлийн сангийн схем (ER Диаграмм)}

Cloud Firestore-ийн гурван collection болон тэдгээрийн хоорондын холбоосыг дараахи диаграммаар илэрхийлэв.

\begin{figure}[h]
\centering
\begin{tikzpicture}[
  coll/.style={draw, rectangle, rounded corners=4pt, fill=green!8,
               minimum width=3.8cm, font=\small\bfseries, text centered},
  field/.style={draw=none, font=\scriptsize, text width=3.6cm, text centered},
  rel/.style={-latex', thick, dashed, gray}
]
% users collection
\node[coll, minimum height=0.7cm] (UT) at (0,9) {users};
\node[field] at (0,8.35) {uid (document ID)};
\node[field] at (0,7.85) {name : string};
\node[field] at (0,7.35) {email : string};
\node[field] at (0,6.85) {emailVerified : boolean};
\node[field] at (0,6.35) {createdAt : timestamp};
\draw[draw] (-1.9,6.05) rectangle (1.9,9.35);

% bookings collection
\node[coll, minimum height=0.7cm] (BT) at (6,9) {bookings};
\node[field] at (6,8.35) {bookingCode (DBK-xxxxx)};
\node[field] at (6,7.85) {userId : string};
\node[field] at (6,7.35) {userEmail : string};
\node[field] at (6,6.85) {userName : string};
\node[field] at (6,6.35) {venueId : string};
\node[field] at (6,5.85) {venueName : string};
\node[field] at (6,5.35) {date : string (YYYY-MM-DD)};
\node[field] at (6,4.85) {timeSlot : string};
\node[field] at (6,4.35) {courtType : full|half};
\node[field] at (6,3.85) {status : active|cancelled|completed};
\node[field] at (6,3.35) {createdAt : timestamp};
\draw[draw] (4.1,3.05) rectangle (7.9,9.35);

% fixed_bookings collection
\node[coll, minimum height=0.7cm] (FT) at (0,4) {fixed\_bookings};
\node[field] at (0,3.35) {venueId : string};
\node[field] at (0,2.85) {organizationName : string};
\node[field] at (0,2.35) {dayOfWeek : number (0--6)};
\node[field] at (0,1.85) {timeSlot : string};
\node[field] at (0,1.35) {isActive : boolean};
\draw[draw] (-1.9,1.05) rectangle (1.9,4.35);

% email_verifications collection
\node[coll, minimum height=0.7cm] (EV) at (6,4) {email\_verifications};
\node[field] at (6,3.35) {uid (document ID)};
\node[field] at (6,2.85) {code : string (6 оронтой)};
\node[field] at (6,2.35) {email : string};
\node[field] at (6,1.85) {expiresAt : timestamp};
\node[field] at (6,1.35) {createdAt : timestamp};
\draw[draw] (4.1,1.05) rectangle (7.9,4.35);

% Relationships
\draw[rel] (1.9,7.85) -- node[above,font=\scriptsize]{userId} (4.1,7.85);
\draw[rel] (1.9,7.35) -- node[above,font=\scriptsize]{userEmail} (4.1,7.35);
\draw[rel] (1.9,6.35) -- node[above,font=\scriptsize]{venueId} (4.1,6.35);
\draw[rel] (1.9,6.85) -- node[below,font=\scriptsize]{uid} (4.1,3.35);
\end{tikzpicture}
\caption{Firestore өгөгдлийн сангийн схем}
\label{fig:erd}
\end{figure}

%-------------------------------------------------------------------------------

\section{Системийн архитектур}

\subsection{Ерөнхий архитектур}

Систем нь \textbf{гурван давхар архитектур}-т тулгуурлана: Flutter мобайл клиент, Node.js/Express бэкэнд (AWS EC2), Firebase үйлчилгээнүүд. Flutter апп нь \code{lib/config.dart}-д тодорхойлсон URL-аар AWS EC2 бэкэндтэй REST API дамжуулан, Firebase SDK-р Cloud Firestore болон Authentication-тай шууд холбогддог.

\begin{figure}[h]
\centering
\begin{tikzpicture}[node distance=2.2cm, auto,
    mobile/.style={rectangle, draw, fill=yellow!15, text width=5.5em,
        text centered, rounded corners, minimum height=3.5em, font=\small},
    backend/.style={rectangle, draw, fill=orange!20, text width=6em,
        text centered, rounded corners, minimum height=3.5em, font=\small},
    cloud/.style={rectangle, draw, fill=blue!15, text width=6em,
        text centered, rounded corners, minimum height=3.5em, font=\small},
    ai/.style={rectangle, draw, fill=green!15, text width=6em,
        text centered, rounded corners, minimum height=3.5em, font=\small},
    line/.style={draw, -latex', thick}]

    % ====== Client Layer ======
    \node [mobile] (app) at (0,0) {Flutter\\Мобайл Апп\\(Android / iOS)};

    % ====== Cloud Backend Layer ======
    \node [backend] (ec2) at (6,0)
        {Node.js/Express\\Docker Container\\AWS EC2 (ap-southeast-1)};

    % ====== Firebase Layer ======
    \node [cloud] (firebase) at (0,-4)
        {Google Firebase\\Auth · Firestore\\Cloud Functions};

    % ====== AI Services ======
    \node [ai] (groq)   at (11,1.5)  {Groq API\\Llama 3.1-8b\\(нөөц горим)};
    \node [ai] (claude) at (11,-1.5) {Anthropic API\\Claude Haiku 4.5\\(үндсэн горим)};

    % ====== DevOps Layer ======
    \node [backend] (ecr) at (6,-4)
        {AWS ECR\\Docker Registry\\(зураг хадгалалт)};
    \node [backend] (gh) at (0,-7)
        {GitHub Repository\\Actions CI/CD};

    % ====== Lines ======
    \path [line] (app)      -- node[above, font=\tiny]{REST API / HTTPS}   (ec2);
    \path [line] (app)      -- node[left,  font=\tiny]{Firebase SDK}        (firebase);
    \path [line] (ec2)      -- node[above, font=\tiny]{Groq API}            (groq);
    \path [line] (ec2)      -- node[below, font=\tiny]{Anthropic API}       (claude);
    \path [line] (firebase) -- node[above, font=\tiny]{Cloud Functions}     (ec2);
    \path [line] (gh)       -- node[below, font=\tiny]{docker push}         (ecr);
    \path [line] (ecr)      -- node[right, font=\tiny]{docker pull \& run}  (ec2);
    \path [line] (gh)       -- node[above, font=\tiny]{git push trigger}    (firebase);
\end{tikzpicture}
\caption{Системийн гурван давхаргат үүлэн архитектурын ерөнхий диаграмм}
\label{fig:arch}
\end{figure}

%-------------------------------------------------

\subsection{Үүлэн байршуулалтын (Cloud Deployment) архитектур}

Систем нь \textbf{AWS үүлэн дэд бүтэц} дээр ажилладаг бөгөөд бэкэнд сервер нь \textbf{Docker контейнер}-ээр \textbf{AWS EC2} дээр байрлана. \textbf{Firebase} нь Google-ийн дэлхийн дэд бүтцэд, \textbf{AI үйлчилгээнүүд} нь тус тусын API үүлэн серверт ажиллана. Апп нь интернет холболттой бол дэлхийн аль ч газраас хандах боломжтой бөгөөд бүх өгөгдөл шифрлэгдэн дамждаг.

\begin{figure}[h]
\centering
\begin{tikzpicture}[
  box/.style={draw, rounded corners=6pt, minimum width=3.5cm,
              minimum height=1.1cm, align=center, font=\small, thick},
  lbox/.style={draw, rounded corners=6pt, minimum width=12cm,
               minimum height=1.6cm, dashed, font=\small\bfseries},
  arr/.style={-latex', thick, blue!60},
  font=\small
]

% ======= Client Zone =======
\draw[dashed, rounded corners=8pt, fill=yellow!5]
  (-1.5,-0.2) rectangle (3.5,1.8);
\node[font=\tiny\bfseries, above] at (1,1.8) {Клиент (Хэрэглэгчийн төхөөрөмж)};
\node[box, fill=yellow!15] (phone) at (1,0.8) {Flutter App\\Android / iOS};

% ======= AWS Zone =======
\draw[dashed, rounded corners=8pt, fill=orange!5]
  (4.5,-0.2) rectangle (10.5,3.8);
\node[font=\tiny\bfseries, above] at (7.5,3.8) {AWS Cloud (ap-southeast-1)};
\node[box, fill=orange!15] (ec2box)  at (7.5,2.8) {AWS EC2 Instance\\t3.micro};
\node[box, fill=orange!10] (docker)  at (7.5,1.5) {Docker Container\\Node.js/Express :3000};
\node[box, fill=orange!10] (ecrbox)  at (7.5,0.2) {AWS ECR\\Docker Image Registry};

% ======= Firebase Zone =======
\draw[dashed, rounded corners=8pt, fill=blue!5]
  (-1.5,-4.8) rectangle (3.5,-2.2);
\node[font=\tiny\bfseries, above] at (1,-2.2) {Google Firebase Cloud};
\node[box, fill=blue!12] (auth)      at (1,-3.2) {Firebase Auth};
\node[box, fill=blue!12] (firestore) at (1,-4.2) {Cloud Firestore};

% ======= CI/CD Zone =======
\draw[dashed, rounded corners=8pt, fill=green!5]
  (4.5,-4.8) rectangle (10.5,-2.2);
\node[font=\tiny\bfseries, above] at (7.5,-2.2) {DevOps / CI-CD Pipeline};
\node[box, fill=green!12] (github)   at (7.5,-3.2) {GitHub Repository};
\node[box, fill=green!12] (actions)  at (7.5,-4.2) {GitHub Actions Workflow};

% ======= Connections =======
\draw[arr] (phone.east)   -- node[above,font=\tiny]{REST API}    (docker.west |- phone.east);
\draw[arr] (phone.south)  -- node[left, font=\tiny]{SDK}          (auth.north);
\draw[arr] (docker.south) -- node[right,font=\tiny]{pull image}   (ecrbox.north);
\draw[arr] (actions.north)-- node[right,font=\tiny]{docker push}  (ecrbox.south);
\draw[arr] (github.west)  -| node[above,font=\tiny]{trigger}      (actions.north -| github.west);

\end{tikzpicture}
\caption{Үүлэн байршуулалтын архитектур --- AWS, Firebase, CI/CD}
\label{fig:deployment}
\end{figure}

%-------------------------------------------------

\subsection{CI/CD Хоолойн архитектур (GitHub Actions)}

\textbf{GitHub Actions}-ын CI/CD хоолой нь хөгжүүлэгч кодоо push хийх бүрт автоматаар ажилладаг. Нийт байршуулалт \textbf{3--4 минут} болдог.

\begin{figure}[h]
\centering
\begin{tikzpicture}[
  cicdstep/.style={draw, rounded corners=5pt, fill=green!8, minimum width=3cm,
               minimum height=0.9cm, align=center, font=\small, thick},
  dec/.style={draw, diamond, fill=orange!15, aspect=2.2,
              align=center, font=\scriptsize, inner sep=2pt},
  arr/.style={-latex', thick},
  font=\small
]

\node[cicdstep] (push)    at (0,8)    {git push\\backend/};
\node[cicdstep] (trigger) at (0,6.8)  {GH Actions\\Trigger};
\node[cicdstep] (build)   at (0,5.6)  {Docker Build\\(Node.js image)};
\node[dec]      (ok)      at (0,4.2)  {Build\\амжилттай?};
\node[cicdstep] (ecr)     at (0,3.0)  {ECR Push\\(AWS Registry)};
\node[cicdstep] (ssh)     at (0,1.8)  {SSH to EC2\\(ap-southeast-1)};
\node[cicdstep] (pull)    at (0,0.6)  {docker pull\\шинэ image};
\node[cicdstep] (run)     at (0,-0.6) {docker stop old\\docker run new};
\node[cicdstep, fill=red!8] (fail) at (4,4.2)  {Workflow fail\\мэдэгдэл};

\draw[arr] (push)    -- (trigger);
\draw[arr] (trigger) -- (build);
\draw[arr] (build)   -- (ok);
\draw[arr] (ok)      -- node[left,font=\scriptsize]{Тийм}  (ecr);
\draw[arr] (ok)      -- node[above,font=\scriptsize]{Үгүй} (fail);
\draw[arr] (ecr)     -- (ssh);
\draw[arr] (ssh)     -- (pull);
\draw[arr] (pull)    -- (run);

\end{tikzpicture}
\caption{GitHub Actions CI/CD хоолойн дэлгэрэнгүй диаграмм}
\label{fig:cicd}
\end{figure}

%-------------------------------------------------

\subsection{Flutter апп-ын бүтэн архитектур}

Flutter апп-ын дотоод архитектур нь \textbf{давхаргат бүтэц}-тэй бөгөөд UI Layer, Service Layer, Data Layer гэсэн гурван давхаргаар зохион байгуулагдсан.

\begin{figure}[h]
\centering
\begin{tikzpicture}[
  layer/.style={draw, rounded corners=6pt, minimum width=13cm,
                minimum height=1.2cm, text centered, font=\small\bfseries, thick},
  comp/.style={draw, rounded corners=4pt, minimum width=2.8cm,
               minimum height=0.85cm, align=center, font=\scriptsize, fill=white},
  arr/.style={-latex', thick, gray},
  font=\small
]

% ===== UI Layer =====
\draw[fill=yellow!10, rounded corners=6pt]
  (-6.5,9.5) rectangle (6.5,12.5);
\node[font=\small\bfseries] at (0,12.2) {UI Layer (lib/screens/ + lib/widgets/)};
\node[comp, fill=yellow!20] at (-4.5,10.8) {HomeScreen\\(заалны жагсаалт)};
\node[comp, fill=yellow!20] at (-1.5,10.8) {BookingScreen\\(цаг захиалах)};
\node[comp, fill=yellow!20] at ( 1.5,10.8) {ChatScreen\\(AI чатбот)};
\node[comp, fill=yellow!20] at ( 4.5,10.8) {ProfileScreen\\(хэрэглэгч)};
\node[comp, fill=yellow!30] at (-3,9.9) {MainScreen\\(BottomNavigationBar)};
\node[comp, fill=yellow!30] at ( 3,9.9) {Custom Widgets\\(VenueCard, SlotGrid)};

% ===== State Management Layer =====
\draw[fill=orange!8, rounded corners=6pt]
  (-6.5,8.0) rectangle (6.5,9.3);
\node[font=\small\bfseries] at (0,9.05) {State Management (StatefulWidget + setState)};
\node[comp, fill=orange!15] at (-3,8.4) {AppState\\(root StatefulWidget)};
\node[comp, fill=orange!15] at ( 3,8.4) {Local State\\(screen-level setState)};

% ===== Service Layer =====
\draw[fill=blue!8, rounded corners=6pt]
  (-6.5,5.8) rectangle (6.5,7.8);
\node[font=\small\bfseries] at (0,7.55) {Service Layer (lib/services/)};
\node[comp, fill=blue!12] at (-4.5,6.7) {AuthService\\(Firebase Auth)};
\node[comp, fill=blue!12] at (-1.5,6.7) {BookingService\\(Firestore + API)};
\node[comp, fill=blue!12] at ( 1.5,6.7) {ChatService\\(Claude/Groq)};
\node[comp, fill=blue!12] at ( 4.5,6.7) {NotificationService\\(flutter\_local)};

% ===== Config Layer =====
\draw[fill=green!8, rounded corners=6pt]
  (-6.5,4.5) rectangle (6.5,5.6);
\node[font=\small\bfseries] at (0,5.35) {Config Layer (lib/config.dart + lib/models/)};
\node[comp, fill=green!12] at (-3.5,4.9) {AppConfig\\(API URLs)};
\node[comp, fill=green!12] at (0,4.9) {Booking Model};
\node[comp, fill=green!12] at (3.5,4.9) {AppTheme\\(Design System)};

% ===== External Services =====
\draw[fill=gray!8, rounded corners=6pt]
  (-6.5,2.5) rectangle (6.5,4.3);
\node[font=\small\bfseries] at (0,4.05) {External Cloud Services};
\node[comp, fill=gray!15] at (-4.5,3.2) {Firebase Auth\\(Google Cloud)};
\node[comp, fill=gray!15] at (-1.5,3.2) {Cloud Firestore\\(Google Cloud)};
\node[comp, fill=gray!15] at ( 1.5,3.2) {AWS EC2\\(Node.js API)};
\node[comp, fill=gray!15] at ( 4.5,3.2) {AI APIs\\(Claude/Groq)};

% ===== Arrows between layers =====
\draw[arr] (0,9.3)  -- (0,8.0);
\draw[arr] (0,8.0)  -- (0,7.8);
\draw[arr] (0,5.8)  -- (0,5.6);
\draw[arr] (0,4.5)  -- (0,4.3);

\end{tikzpicture}
\caption{Flutter аппликейшний бүтэн давхаргат архитектур}
\label{fig:flutter_arch}
\end{figure}

Flutter апп-д дараахи хавтасны бүтцийг ашигласан:

\begin{itemize}
    \item \code{lib/screens/} — Дэлгэцийн файлууд (\code{home\_screen.dart}, \code{booking\_screen.dart}, \code{chat\_screen.dart}, \code{profile\_screen.dart});
    \item \code{lib/widgets/} — Дахин ашиглах widget-ууд (\code{venue\_card.dart}, \code{time\_slot\_grid.dart}, \code{booking\_card.dart});
    \item \code{lib/models/} — Өгөгдлийн загвар классууд (\code{Booking}, \code{FixedBooking});
    \item \code{lib/services/} — Firebase болон API үйлчилгээний классууд (\code{auth\_service.dart}, \code{booking\_service.dart}, \code{chat\_service.dart});
    \item \code{lib/config.dart} — API URL болон тохиргооны нэг эх сурвалж (\code{AppConfig}, \code{AppTheme}).
\end{itemize}

Апп-ын навигацийг \code{BottomNavigationBar} widget дамжуулан 4 tab-тай хэрэгжүүлсэн бөгөөд идэвхтэй tab-ийн индексийг root \code{StatefulWidget}-ийн \code{setState}-д хадгалдаг.

\subsection{API URL-ийн нэг эх сурвалж (lib/config.dart)}

Апп-ын бүх API хаягийг \code{lib/config.dart} файлд төвлөрүүлснээр хөгжүүлэлтийн болон продакшн орчинд шилжихэд нэг файлыг л өөрчлөх хэрэгтэй болдог. Систем нь \textbf{AWS үүлэн орчинд} бодитоор ажиллаж байгаа тул продакшн URL нь EC2 instance-ийн нийтийн IP хаяг юм:

\begin{lstlisting}[language={}, caption={lib/config.dart --- AWS үүлэн орчны API URL тохиргоо}]
class AppConfig {
  // Local test: Android emulator -> 'http://10.0.2.2:3000'
  //             iOS simulator    -> 'http://localhost:3000'
  // Production: AWS EC2 (ap-southeast-1, Singapore) - Cloud deployment
  static const String backendUrl = 'http://<EC2-PUBLIC-IP>:3000';

  static const String authEndpoint     = '$backendUrl/api/auth';
  static const String chatEndpoint     = '$backendUrl/api/chat';
  static const String usersEndpoint    = '$backendUrl/api/users';
  static const String bookingsEndpoint = '$backendUrl/api/bookings';
  static const String scheduleEndpoint = '$backendUrl/api/schedule';
}
\end{lstlisting}

%-------------------------------------------------------------------------------

\section{Өгөгдлийн сангийн зохион байгуулалт}

Cloud Firestore-д дараахи гурван collection бүтцийг ашигласан. Спорт заалнуудын мэдээлэл Firestore-д хадгалагддаггүй — кодод шууд тодорхойлсон (hardcoded) байдлаар тавьж, Firestore-ийн уншилтын тоог хэмнэдэг.

\subsection{\texttt{users} collection}

Хэрэглэгчийн мэдээлэл:
\begin{itemize}
    \item \code{uid} (string) — Firebase Auth-ын хэрэглэгчийн ID (document ID болгон ашиглана);
    \item \code{name} (string) — Нэр;
    \item \code{email} (string) — Email хаяг;
    \item \code{emailVerified} (boolean) — Email баталгаажсан эсэх;
    \item \code{createdAt} (timestamp) — Бүртгэсэн огноо.
\end{itemize}

\subsection{\texttt{bookings} collection}

Захиалгын мэдээлэл:
\begin{itemize}
    \item \code{code} (string) — DBK-xxxxx форматын захиалгын код;
    \item \code{userId} (string) — Хэрэглэгчийн Firebase UID;
    \item \code{userEmail} (string) — Хэрэглэгчийн email хаяг;
    \item \code{userName} (string) — Хэрэглэгчийн нэр;
    \item \code{venueId} (string) — Заалны ID (hardcoded жагсаалтаас);
    \item \code{venueName} (string) — Заалны нэр;
    \item \code{venueType} (string) — Заалны спортийн төрөл (сагсан бөмбөг, волейбол гэх мэт);
    \item \code{venueLocation} (string) — Заалны байршил;
    \item \code{date} (Timestamp) — Захиалгын огноо;
    \item \code{dateKey} (string) — YYYY-MM-DD форматын огноо (Firestore шүүлтийн хурдыг нэмэгдүүлэх зорилгоор);
    \item \code{timeSlot} (string) — Цагийн эхлэл (жишээ: "09:00");
    \item \code{timeSlotEnd} (string) — Цагийн төгсгөл (жишээ: "10:00");
    \item \code{courtType} (string) — Талбайн төрөл: \code{full} (бүтэн талбай) | \code{half} (хагас талбай);
    \item \code{price} (string) — Захиалгын үнэ;
    \item \code{status} (string) — Төлөв: \code{upcoming} | \code{cancelled} | \code{completed};
    \item \code{createdAt} (timestamp) — Үүсгэсэн огноо.
\end{itemize}

\subsection{\texttt{fixed\_bookings} collection}

Байгууллагуудын тогтмол давтагдах захиалга:
\begin{itemize}
    \item \code{venueId} (string) — Заалны ID;
    \item \code{organizationName} (string) — Байгууллагын нэр (засгийн газар, компани);
    \item \code{dayOfWeek} (number) — Долоо хоногийн өдөр (0=Ням, 1=Даваа \ldots 6=Бямба);
    \item \code{timeSlot} (string) — Хаагдсан цагийн үе;
    \item \code{isActive} (boolean) — Идэвхтэй эсэх.
\end{itemize}

Тогтмол захиалгыг хэрэглэгчийн цагийн хуваарь харах үед уншиж, тухайн өдрийн хаагдсан цагийн үеийг харуулна. Ингэснээр засгийн газар, компаниудын долоо хоногийн тогтмол цагийг захиалгаас урьдчилан хасдаг.

\subsection{\texttt{email\_verifications} collection}

Бүртгэлийн баталгаажуулах кодыг түр хадгалах зориулалттай collection:
\begin{itemize}
    \item Document ID — Firebase Auth-ийн хэрэглэгчийн UID;
    \item \code{code} (string) — 6 оронтой санамсаргүй тоон код;
    \item \code{email} (string) — Кодыг илгээсэн email хаяг;
    \item \code{expiresAt} (timestamp) — Кодын хугацаа (үүсгэснээс 10 минутын дараа);
    \item \code{createdAt} (timestamp) — Үүсгэсэн огноо.
\end{itemize}

Баталгаажуулалт амжилттай болсны дараа тухайн document автоматаар устгагддаг. Хугацаа хэтэрсэн код дахин илгээх шаардлагатай болно.

\subsection{Hardcoded спорт заалнуудын жагсаалт}

Таван спорт заалны мэдээлэл Firestore collection биш, Flutter кодод шууд тодорхойлсон:

\begin{lstlisting}[language={}, caption={Hardcoded спорт заалнуудын жагсаалт}]
const List<Map<String, dynamic>> venues = [
  {'id': 'goviin_arena',     'name': 'Goviin Arena'},
  {'id': 'omnogovi_sport',   'name': 'Omnogovi Sport Hall'},
  {'id': 'goviin_voleybol',  'name': 'Goviin Volleyball Club'},
  {'id': 'stadion_voleybol', 'name': 'Stadium Volleyball Hall'},
  {'id': 'central_basket',   'name': 'Central Basketball Hall'},
];
\end{lstlisting}

%-------------------------------------------------------------------------------

\section{Flutter апп-ын хэрэгжүүлэлт}

\subsection{Дулаан говийн дизайн (Warm Desert Theme)}

Апп-ын дизайн нь говийн цөлийн дулаан тонийг илэрхийлсэн, нүдэнд ээлтэй цайвар дэвсгэр дээрх наранжин-хүрэн өнгөн схемтэй. \code{Montserrat} фонт ашигласан бөгөөд дөрвөн үндсэн өнгөөр тодорхойлогддог:

\begin{itemize}
    \item \textbf{Дулаан бежийн дэвсгэр} (\code{\#EDE7DA}) — Scaffold арын фон;
    \item \textbf{Наранжин} (\code{\#D4700F}) — Товч, идэвхтэй элемент (secondary);
    \item \textbf{Гүн наранжин} (\code{\#B85C08}) — Дэд онцлох (accent);
    \item \textbf{Харанхуй хүрэн} (\code{\#1A1208}) — Үндсэн текст.
\end{itemize}

Material Design 3-ийн \code{ThemeData} классаар дамжуулан апп-ын бүх widget-д нэгдсэн цайвар дизайн хэрэгжүүлсэн:

\begin{lstlisting}[language={}, caption={Material Design 3 дулаан говийн дизайн тохиргоо}]
class AppTheme {
  static const Color primary       = Color(0xFFEDE7DA); // beige background
  static const Color secondary     = Color(0xFFD4700F); // orange
  static const Color accent        = Color(0xFFB85C08); // deep orange
  static const Color cardColor     = Color(0xFFF8F4EE); // light card
  static const Color textPrimary   = Color(0xFF1A1208); // dark brown

  static ThemeData get lightTheme => ThemeData(
    useMaterial3: true,
    fontFamily: 'Montserrat',
    brightness: Brightness.light,
    scaffoldBackgroundColor: primary,
    colorScheme: const ColorScheme.light(
      primary:   secondary,
      secondary: accent,
      surface:   Color(0xFFE4DDD0),
    ),
  );
}
\end{lstlisting}

\subsection{4 tab-тай доод навигаци}

Апп-ын үндсэн навигаци нь \code{BottomNavigationBar} ашиглан 4 tab-тай:

\begin{lstlisting}[language={}, caption={BottomNavigationBar — 4 tab навигаци}]
class MainScreen extends StatefulWidget {
  const MainScreen({super.key});
  @override
  State<MainScreen> createState() => _MainScreenState();
}

class _MainScreenState extends State<MainScreen> {
  int _selectedIndex = 0;

  final List<Widget> _screens = const [
    HomeScreen(),
    BookingsScreen(),
    ChatScreen(),
    ProfileScreen(),
  ];

  @override
  Widget build(BuildContext context) {
    return Scaffold(
      body: _screens[_selectedIndex],
      bottomNavigationBar: BottomNavigationBar(
        currentIndex: _selectedIndex,
        onTap: (index) => setState(() => _selectedIndex = index),
        items: const [
          BottomNavigationBarItem(icon: Icon(Icons.home),
              label: 'Home'),
          BottomNavigationBarItem(icon: Icon(Icons.calendar_today),
              label: 'Bookings'),
          BottomNavigationBarItem(icon: Icon(Icons.chat),
              label: 'Chat'),
          BottomNavigationBarItem(icon: Icon(Icons.person),
              label: 'Profile'),
        ],
      ),
    );
  }
}
\end{lstlisting}

\subsection{Email/нууц үгийн нэвтрэлт ба баталгаажуулалт}

Нэвтрэлтийн дэлгэцэд хэрэглэгч \textbf{email хаяг} болон \textbf{нууц үгээр} бүртгүүлж нэвтрэх боломжтой. Анхны бүртгэлийн үед тухайн email хаягт \textbf{6 оронтой баталгаажуулах код} илгээгдэх бөгөөд кодоо зөв оруулсны дараа л нэвтрэх боломж нээгдэнэ.

\textbf{Бүртгэлийн урсгал:}
\begin{enumerate}
    \item Хэрэглэгч нэр, email, нууц үгээ оруулна;
    \item Firebase Auth-д \code{createUserWithEmailAndPassword} дуудаж UID үүсгэж, Firestore \code{users} collection-д \code{emailVerified: false} тэмдэглэн хадгална;
    \item Backend-ийн \code{POST /api/auth/send-code} endpoint-д UID болон email илгээнэ;
    \item Backend нь Nodemailer-ээр 6 оронтой код generate хийж email рүү илгээнэ; Firestore \code{email\_verifications} collection-д 10 минутын хугацаатайгаар хадгална;
    \item Хэрэв код илгээхэд алдаа гарвал Firebase Auth хэрэглэгч болон Firestore бүртгэлийг автоматаар устгаж (rollback), хэрэглэгч дахин оролдох боломж нээгддэг;
    \item Хэрэглэгч кодоо оруулахад \code{POST /api/auth/verify-code} дуудаж баталгаажуулна;
    \item Баталгаажсан тохиолдолд Firestore-д \code{emailVerified: true} тэмдэглэгдэж, \code{email\_verifications} document устгагдан, апп нүүр дэлгэц рүү шилжинэ.
\end{enumerate}

\begin{lstlisting}[language={}, caption={Firebase Email Authentication — бүртгэл}]
Future<void> registerWithEmail(
    String email, String password, String name) async {
  final credential = await FirebaseAuth.instance
      .createUserWithEmailAndPassword(
        email: email, password: password);
  final uid = credential.user!.uid;

  // Email verification code via backend
  await http.post(
    Uri.parse(AppConfig.authEndpoint + '/send-code'),
    headers: {'Content-Type': 'application/json'},
    body: jsonEncode({'email': email, 'uid': uid}),
  );
}
\end{lstlisting}

\begin{lstlisting}[language={}, caption={Backend — Nodemailer ашиглан код илгээх}]
// POST /api/auth/send-code
router.post('/send-code', async (req, res) => {
  const { email, uid } = req.body;
  const code = Math.floor(100000
      + Math.random() * 900000).toString();

  await db.collection('email_verifications').doc(uid).set({
    code, email,
    expiresAt: new Date(Date.now() + 10 * 60 * 1000),
  });

  const transporter = nodemailer.createTransport({
    host: 'smtp.gmail.com', port: 587,
    auth: { user: process.env.EMAIL_USER,
            pass: process.env.EMAIL_PASS },
  });

  await transporter.sendMail({
    from: `"Goviin Sport" <${process.env.EMAIL_USER}>`,
    to: email,
    subject: 'Goviin Sport - Batalgaajuulakh kod',
    html: `<h2>${code}</h2><p>10 minutiin dotoor huchintei</p>`,
  });

  res.json({ success: true });
});
\end{lstlisting}

Баталгаажуулалт амжилттай болсны дараа Firebase нь хэрэглэгчийн UID үүсгэж, Firestore-ийн \code{users} collection-д email болон \code{emailVerified: true} талбартайгаар бүртгэл нэмнэ. Нэвтрэх үед backend \code{emailVerified} талбарыг шалгаж, баталгаажаагүй бол нэвтрэхийг хориглоно.

\begin{figure}[h]
\centering
\begin{minipage}[t]{0.30\textwidth}
  
\includegraphics[width=\linewidth]{screenshot_login.png}
  \subcaption{Нэвтрэх дэлгэц}
\end{minipage}
\hfill
\begin{minipage}[t]{0.30\textwidth}
  
\includegraphics[width=\linewidth]{screenshot_register.png}
  \subcaption{Бүртгүүлэх дэлгэц}
\end{minipage}
\hfill
\begin{minipage}[t]{0.30\textwidth}
  \includegraphics[width=\linewidth]{screenshot_verify.png}
  \subcaption{Email баталгаажуулах}
\end{minipage}
\caption{Email/нууц үгийн нэвтрэлт ба баталгаажуулалтын дэлгэцүүд}
\label{fig:auth_screens}
\end{figure}

\subsection{Нүүр дэлгэц — Спорт заалнуудын жагсаалт}

Нүүр дэлгэцэд таван спорт заалын карт харагдана. Хэрэглэгч карт дарахад тухайн заалны хуваарийн дэлгэц нээгдэнэ. Заалнуудын мэдээлэл кодод тодорхойлсон тул Firestore хандалт шаардлагагүй — дэлгэц нэн даруй ачаалдаг.

\begin{figure}[h]
\centering
\begin{minipage}[t]{0.45\textwidth}
  \includegraphics[width=\linewidth]{screenshot_home.png}
  \subcaption{Нүүр дэлгэц — заалны жагсаалт ба шүүлтүүр}
\end{minipage}
\hfill
\begin{minipage}[t]{0.45\textwidth}
  
\includegraphics[width=\linewidth]{screenshot_bookings.png}
  \subcaption{Миний захиалгууд}
\end{minipage}
\caption{Нүүр дэлгэц ба захиалгын жагсаалт}
\label{fig:home_screens}
\end{figure}

\subsection{Цагийн хуваарь ба захиалгын дэлгэц}

Цагийн хуваарийн дэлгэцэд хэрэглэгч огноо сонгохоор тухайн өдрийн цагийн үе бүр харагдана. Нийт \textbf{16 цагийн үе} (08:00–24:00, нэг цагийн алхамтай, шөнийн 12 цаг хүртэл) байдаг. Firestore-аас уг заал, уг өдрийн \code{bookings} болон \code{fixed\_bookings} өгөгдлийг уншиж, захиалагдсан цагийн үеийг идэвхгүй байдлаар харуулна:

\begin{lstlisting}[language={}, caption={Цагийн хуваарь уншиж харуулах}]
Future<void> _loadSlots(String venueId, String date) async {
  setState(() => isLoading = true);

  // User bookings
  final bookSnap = await FirebaseFirestore.instance
      .collection('bookings')
      .where('venueId', isEqualTo: venueId)
      .where('date', isEqualTo: date)
      .where('status', whereIn: ['upcoming', 'active'])
      .get();

  // Fixed/organizational bookings
  final weekday = DateTime.parse(date).weekday % 7;
  final fixedSnap = await FirebaseFirestore.instance
      .collection('fixed_bookings')
      .where('venueId', isEqualTo: venueId)
      .where('dayOfWeek', isEqualTo: weekday)
      .where('isActive', isEqualTo: true)
      .get();

  final bookedSlots = {
    ...bookSnap.docs.map((d) => d['timeSlot'] as String),
    ...fixedSnap.docs.map((d) => d['timeSlot'] as String),
  };

  setState(() {
    _bookedSlots = bookedSlots;
    isLoading = false;
  });
}
\end{lstlisting}

Захиалга хийхдээ \textbf{DBK-xxxxx} форматын захиалгын код автоматаар үүснэ:

\begin{lstlisting}[language={}, caption={Захиалгын код үүсгэх ба Firestore-д хадгалах}]
String _generateBookingCode() {
  final rand = Random();
  final digits = List.generate(5,
      (_) => rand.nextInt(10).toString()).join();
  return 'DBK-$digits';
}

Future<String> saveBooking({...}) async {
  final code = 'DBK-${DateTime.now().millisecondsSinceEpoch % 10000}';
  final dateKey = '${date.year}-${date.month.toString()
      .padLeft(2,'0')}-${date.day.toString().padLeft(2,'0')}';

  final doc = await _bookings.add({
    'venueId':        venueId,
    'venueName':      venueName,
    'venueType':      venueType,
    'venueLocation':  venueLocation,
    'date':           Timestamp.fromDate(date),
    'dateKey':        dateKey,        // fast filter key
    'timeSlot':       timeSlot,
    'timeSlotEnd':    timeSlotEnd,
    'courtType':      courtType,      // full | half
    'price':          price,
    'userName':       userName,
    'userEmail':      userEmail,
    'status':         'upcoming',     // not active yet
    'code':           code,
    'createdAt':      FieldValue.serverTimestamp(),
  });
  return doc.id;
}
\end{lstlisting}

\begin{figure}[h]
\centering
\begin{minipage}[t]{0.45\textwidth}
  
\includegraphics[width=\linewidth]{screenshot_booking.png}
  \subcaption{Цаг ба талбайн төрөл сонгох дэлгэц}
\end{minipage}
\hfill
\begin{minipage}[t]{0.45\textwidth}
  
\includegraphics[width=\linewidth]{screenshot_confirm.png}
  \subcaption{Захиалга баталгаажуулах дэлгэц}
\end{minipage}
\caption{Цагийн хуваарь ба захиалгын дэлгэцүүд}
\label{fig:booking_screens}
\end{figure}

\subsection{Локал мэдэгдлийн хэрэгжүүлэлт}

\code{flutter\_local\_notifications} ашиглан захиалгаас 1 цагийн өмнө сануулга schedule хийдэг. FCM сервер байхгүй тул энэ нь бүрэн офлайн ажилладаг:

\begin{lstlisting}[language={}, caption={flutter\_local\_notifications — захиалгын сануулга}]
Future<void> _scheduleReminder(
    String timeSlot, String bookingCode) async {
  final parts = timeSlot.split(':');
  final slotHour = int.parse(parts[0]);
  final slotMin  = int.parse(parts[1].split('-')[0]);

  final bookingTime = DateTime(
    DateTime.now().year, DateTime.now().month,
    DateTime.now().day, slotHour, slotMin,
  );
  final reminderTime = bookingTime.subtract(
      const Duration(hours: 1));

  await flutterLocalNotificationsPlugin.zonedSchedule(
    bookingCode.hashCode,
    'Booking Reminder',
    '$bookingCode - starts in 1 hour',
    tz.TZDateTime.from(reminderTime, tz.local),
    notificationDetails,
    androidScheduleMode:
        AndroidScheduleMode.exactAllowWhileIdle,
    uiLocalNotificationDateInterpretation:
        UILocalNotificationDateInterpretation
            .absoluteTime,
  );
}
\end{lstlisting}

\begin{figure}[h]
\centering
\includegraphics[width=0.35\linewidth]{screenshot_notifications.png}
\caption{Мэдэгдлүүдийн дэлгэц --- захиалгын сануулга}
\label{fig:notifications_screen}
\end{figure}

%-------------------------------------------------------------------------------

\section{Node.js/Express бэкэндийн хэрэгжүүлэлт}

\subsection{Бэкэндийн бүтэц}

Node.js/Express бэкэнд нь \code{backend/} хавтаст байрлах бөгөөд дараахи бүтэцтэй:

\begin{itemize}
    \item \code{server.js} — Үндсэн Express сервер, middleware, route бүртгэл;
    \item \code{routes/auth.js} — Email баталгаажуулалтын endpoint (\code{/api/auth});
    \item \code{routes/bookings.js} — Захиалгын удирдлага, cron ажлын endpoint;
    \item \code{routes/users.js} — Хэрэглэгчийн мэдээлэл;
    \item \code{routes/schedule.js} — Цагийн хуваарийн мэдээлэл.
\end{itemize}

\subsection{Үндсэн API endpoint-ууд}

\begin{lstlisting}[language={}, caption={Express.js — үндсэн route тохиргоо}]
const express = require('express');
const app     = express();

app.use(express.json());
app.use('/api/auth',     require('./routes/auth'));
app.use('/api/chat',     require('./routes/chat'));
app.use('/api/users',    require('./routes/users'));
app.use('/api/bookings', require('./routes/bookings'));
app.use('/api/schedule', require('./routes/schedule'));

app.listen(3000, () => {
  console.log('Server running on port 3000');
});
\end{lstlisting}

\subsection{Cron ажил — Дуусгавар болсон захиалгыг автоматаар шинэчлэх}

Бэкэнд нь 15 минут тутамд cron ажлыг ажиллуулж, эхлэх цаг нь өнгөрсөн боловч \code{active} төлөвтэй хэвээр байгаа захиалгыг \code{completed} болгон шинэчилдэг. Энэ нь Firestore-д шууд хандан batch update хийдэг:

\begin{lstlisting}[language={}, caption={Cron ажил — 15 минут тутамд захиалга шинэчлэх}]
const cron = require('node-cron');
const admin = require('firebase-admin');

// runs every 15 minutes
cron.schedule('*/15 * * * *', async () => {
  const now   = new Date();
  const today = now.toISOString().split('T')[0];
  const nowTime = now.getHours().toString().padStart(2,'0')
                + ':' + now.getMinutes().toString().padStart(2,'0');

  const snapshot = await admin.firestore()
    .collection('bookings')
    .where('status', '==', 'active')
    .where('date', '<=', today)
    .get();

  const batch = admin.firestore().batch();
  snapshot.docs.forEach(doc => {
    const endTime = doc.data().timeSlot.split('-')[1];
    if (doc.data().date < today ||
        (doc.data().date === today && endTime <= nowTime)) {
      batch.update(doc.ref, { status: 'completed' });
    }
  });
  await batch.commit();
  console.log('Cron: expired bookings completed');
});
\end{lstlisting}

%-------------------------------------------------------------------------------

\section{AI чатботын хэрэгжүүлэлт}

\subsection{Архитектур — Үндсэн болон нөөц горим}

AI чатботын архитектур нь хоёр давхар горимтой:

\begin{enumerate}
    \item \textbf{Үндсэн горим} — Firebase Cloud Functions дамжуулан Anthropic Claude Haiku 4.5;
    \item \textbf{Нөөц горим} — Node.js Express бэкэнд дамжуулан Groq Llama 3.1-8b.
\end{enumerate}

Flutter апп нь \code{AppConfig.chatEndpoint} (Express бэкэнд) болон Firebase Cloud Function-ийн URL аль алинтай холбогдох чадвартай.

\subsection{Firebase Cloud Functions — Claude Haiku}

Cloud Functions дотор Anthropic-ийн Node.js SDK ашиглан Claude Haiku 4.5 API-тай харилцана:

\begin{lstlisting}[language={}, caption={Firebase Cloud Functions — Claude Haiku чат}]
const functions = require('firebase-functions');
const Anthropic = require('@anthropic-ai/sdk');

exports.chat = functions.https.onCall(async (data, context) => {
  if (!context.auth) {
    throw new functions.https.HttpsError('unauthenticated',
        'Authentication required');
  }

  const client = new Anthropic({
    apiKey: functions.config().anthropic.key,
  });

  const message = await client.messages.create({
    model: 'claude-haiku-4-5',
    max_tokens: 512,
    messages: [{ role: 'user', content: data.message }],
    system: 'You are a sports hall booking assistant for Dalanzadgad city.',
  });

  return { reply: message.content[0].text };
});
\end{lstlisting}

\subsection{Express бэкэнд — Groq Llama нөөц горим}

\begin{lstlisting}[language={}, caption={Express бэкэнд — Groq Llama чат endpoint}]
const Groq = require('groq-sdk');
const router = require('express').Router();

const groq = new Groq({
  apiKey: process.env.GROQ_API_KEY
});

router.post('/', async (req, res) => {
  const { message } = req.body;
  const completion = await groq.chat.completions.create({
    model: 'llama-3.1-8b-instant',
    messages: [
      { role: 'system',
        content: 'You are a sports hall booking assistant.' },
      { role: 'user', content: message },
    ],
    max_tokens: 512,
  });
  res.json({
    reply: completion.choices[0].message.content
  });
});

module.exports = router;
\end{lstlisting}

\begin{figure}[h]
\centering

\includegraphics[width=0.35\linewidth]{screenshot_chat.png}
\caption{AI Туслах дэлгэц --- спорт заалтай холбоотой асуулт асуух}
\label{fig:chat_screen}
\end{figure}

%-------------------------------------------------------------------------------

\section{Docker контейнержуулалт ба AWS байршуулалт}

\subsection{Docker Image}

Node.js бэкэнд нь \code{backend/Dockerfile}-аар тодорхойлогдсон Docker дүрс ашигладаг. Node 18 Alpine суурь дүрс сонгосноор эцсийн дүрсний хэмжээ жижиг байна:

\begin{lstlisting}[language={}, caption={backend/Dockerfile}]
FROM node:18-alpine
WORKDIR /app
COPY package*.json ./
RUN npm ci --only=production
COPY . .
EXPOSE 3000
CMD ["node", "server.js"]
\end{lstlisting}

\subsection{GitHub Actions CI/CD Workflow}

\code{.github/workflows/deploy.yml} файл нь \code{backend/} хавтаст push хийх үед автоматаар ажилладаг:

\begin{lstlisting}[language={}, caption={GitHub Actions — CI/CD workflow}]
name: Deploy Backend

on:
  push:
    paths: ['backend/**']
    branches: [main]

jobs:
  deploy:
    runs-on: ubuntu-latest
    steps:
      - uses: actions/checkout@v3

      - name: Configure AWS credentials
        uses: aws-actions/configure-aws-credentials@v4
        with:
          aws-access-key-id:     ${{ secrets.AWS_ACCESS_KEY_ID }}
          aws-secret-access-key: ${{ secrets.AWS_SECRET_ACCESS_KEY }}
          aws-region: ap-southeast-1

      - name: Login to ECR
        id: login-ecr
        uses: aws-actions/amazon-ecr-login@v2

      - name: Build and push Docker image
        env:
          ECR_REGISTRY: ${{ steps.login-ecr.outputs.registry }}
        run: |
          docker build -t $ECR_REGISTRY/dz-zaal-backend:latest ./backend
          docker push $ECR_REGISTRY/dz-zaal-backend:latest

      - name: Deploy to EC2
        env:
          EC2_HOST: ${{ secrets.EC2_HOST }}
          EC2_USER: ${{ secrets.EC2_USER }}
          EC2_SSH_KEY: ${{ secrets.EC2_SSH_KEY }}
          ECR_REGISTRY: ${{ steps.login-ecr.outputs.registry }}
          EMAIL_USER: ${{ secrets.EMAIL_USER }}
          EMAIL_PASS: ${{ secrets.EMAIL_PASS }}
        run: |
          printf '%s\n' "$EC2_SSH_KEY" | tr -d '\r' > /tmp/deploy_key
          chmod 600 /tmp/deploy_key
          ssh -i /tmp/deploy_key -o StrictHostKeyChecking=no \
            ${EC2_USER:-ec2-user}@$EC2_HOST "
              docker pull $ECR_REGISTRY/dz-zaal-backend:latest
              docker stop dz-zaal-backend || true
              docker rm   dz-zaal-backend || true
              docker run -d --name dz-zaal-backend \
                --restart unless-stopped \
                --env-file ~/backend/.env \
                -e EMAIL_USER=$EMAIL_USER \
                -e EMAIL_PASS=$EMAIL_PASS \
                -p 3000:3000 \
                $ECR_REGISTRY/dz-zaal-backend:latest
            "
\end{lstlisting}

CI/CD хоолойн тусламжтайгаар хөгжүүлэгч \code{git push} хийхэд л бэкэнд автоматаар шинэчлэгддэг. Хүний гараар сервер дээр нэвтрэх шаардлага байхгүй болдог.

%-------------------------------------------------------------------------------

\section{Туршилт ба үр дүнгийн үнэлгээ}

\subsection{Функциональ туршилт}

Апп-ын үндсэн функцүүдийг дараахи тест тохиолдлуудаар шалгасан:

\begin{table}[h]
\centering
\caption{Функциональ туршилтын үр дүн}
\begin{tabular}{|c|p{5.5cm}|c|}
\hline
\textbf{№} & \textbf{Тест тохиолдол} & \textbf{Үр дүн} \\ \hline
1 & Email/нууц үгээр бүртгүүлэх, email баталгаажуулах код амжилттай ирэх & Тэнцсэн \\ \hline
2 & Нүүр дэлгэцэд 5 спорт заал гарах & Тэнцсэн \\ \hline
3 & fixed\_bookings тухайн өдрийн хуваарьт тусах & Тэнцсэн \\ \hline
4 & DBK-xxxxx форматын захиалгын код үүсэх & Тэнцсэн \\ \hline
5 & Захиалга хийснээс 1 цагийн өмнө локал мэдэгдэл ирэх & Тэнцсэн \\ \hline
6 & Захиалга цуцлах, цаг чөлөөлөгдөх & Тэнцсэн \\ \hline
7 & 15 минутын cron ажил дуусгавар болсон захиалгыг completed болгох & Тэнцсэн \\ \hline
8 & AI чат (Claude Haiku) хариулт өгөх & Тэнцсэн \\ \hline
9 & AI чат (Groq Llama нөөц) хариулт өгөх & Тэнцсэн \\ \hline
10 & GitHub Actions push хийснээс хойш EC2 шинэчлэгдэх & Тэнцсэн \\ \hline
\end{tabular}
\label{tab:test}
\end{table}

\subsection{Гүйцэтгэлийн хэмжилт}

Апп-ын гүйцэтгэлийг дараахи үзүүлэлтээр хэмжсэн:

\begin{itemize}
    \item \textbf{Flutter апп ачаалах хугацаа} — Дундаж 2.3 секунд (зорилт: 3 секундаас бага);
    \item \textbf{Firestore хуваарь унших хугацаа} — Дундаж 0.6 секунд;
    \item \textbf{Node.js API хариулах хугацаа (EC2)} — Дундаж 0.8 секунд;
    \item \textbf{Claude Haiku хариулт (Cloud Functions)} — Дундаж 1.4 секунд;
    \item \textbf{Groq Llama хариулт (EC2)} — Дундаж 0.9 секунд;
    \item \textbf{Дэлгэц хоорондох шилжилт} — 60 FPS-д жигд ажиллах.
\end{itemize}

\subsection{CI/CD хоолойн туршилт}

GitHub Actions workflow-ийг дараахи тест тохиолдлуудаар шалгасан:

\begin{itemize}
    \item \code{backend/} хавтаст өөрчлөлт push хийхэд workflow автоматаар trigger болдог болохыг баталгаажуулсан;
    \item Docker дүрс AWS ECR-д амжилттай push хийгдэж, SHA тагаар хянагдаж байгааг баталгаажуулсан;
    \item EC2-д SSH холболтоор шинэ контейнер ажиллаж, хуучин контейнер зогсдог болохыг шалгасан;
    \item Нийт байршуулалтын хугацаа 3--4 минут болохыг хэмжсэн.
\end{itemize}
